\documentclass[12pt]{article}

\input{../../_assets/preambulo.tex}


\begin{document}

    % 1. Foto de fondo
    % 2. Título
    % 3. Encabezado Izquierdo
    % 4. Color de fondo
    % 5. Coord x del titulo
    % 6. Coord y del titulo
    % 7. Fecha

    
    \input{../../_assets/portada}
    \portadaExamen{ffccA4.jpg}{Curvas y\\Superficies\\Examen V}{Curvas y Superficies. Examen V}{MidnightBlue}{-8}{28}{2026}{José Juan Urrutia Milán}

    \begin{description}
        \item[Asignatura] Curvas y Superficies.
        \item[Curso Académico] 2024/25.
        \item[Grado] Grado en Matemáticas.
        % \item[Grupo] B.
        % \item[Profesor] ---.
        \item[Descripción] Convocatoria Ordinaria.
        \item[Fecha] 17 de junio de 2025.
        % \item[Duración] 2 horas y media.
    \end{description}
    \newpage


    % ------------------------------------

    \begin{ejercicio}
        Para cada $r\in \mathbb{R}\setminus \{-1,1\}$ consideramos la curva $\alpha_r:\mathbb{R}\to\mathbb{R}^2$, $\alpha_r(t) = (rt+\cos(t),\sen(t))$.
        \begin{enumerate}[label=\alph*)]
            \item Probar que $\alpha_r$ es regular y determina curvatura de $\alpha_r$.
            \item Deducir que $k_{\alpha_r}$ no se anula si y solo si $|r| < 1$, y que en ese caso $k_{\alpha_r} \geq \frac{1-|r|}{{(|r|+1)}^{3}}$.
        \end{enumerate}
    \end{ejercicio}

    \begin{ejercicio}
        Consideremos el cono $C = \{(x,y,z)\in \mathbb{R}^3 : 3(x^2+y^2)=z^2, z>0\}$ y la aplicación
        \Func{F}{\bb{R}^+\times \bb{R}}{\bb{R}^3}{(r,\theta)}{r(\cos\theta,\sen\theta,\sqrt{3})}
        \begin{enumerate}[label=\alph*)]
            \item Probar que $C$ es una superficie regular orientable y $F(\left]0,+\infty\right[\times \mathbb{R}) = C$. Demostrar que para cada $s\in \mathbb{R}$ la aplicación $X_s:\mathbb{R}^+\times \left]s,s+2\pi\right[\to \mathbb{R}^3$, $X = F\big|_{\mathbb{R}^+\times \left]s,s+2\pi\right[}$ es una parametrización de $C$.
            \item Demostrar que existe una única aplicación de Gauss $N:C\to \bb{S}^2$ tal que $N\circ F = \frac{F_r\times F_\theta}{\|F_r\times F_\theta\|}$. Determinar la expresión local en la carta $X_s$ de $C,s\in \mathbb{R}$, de la primrea forma fundamental, de la segunda forma fundamental, el endomorfismo de Weingarten, la curvatura de Gauss, la curvatura media y las curvaturas y direcciones principales (en la dirección de $N$).
            \item Identificando $\mathbb{R}^2\setminus \{(0,0)\}$ con la superficie regular $(\mathbb{R}^2\setminus \{(0,0)\})\times \{0\}$ de $\mathbb{R}^3$, probar que la aplicación
                \Func{G}{\bb{R}^2\setminus \{(0,0)\}}{C}{(u,v)}{\frac{1}{2\sqrt{u^2+v^2}}\left(u^2-v^2,2uv,\sqrt{3}(u^2+v^2)\right)}
                es una isometría local entre superficies regulares. Determinar para cada $p\in \mathbb{R}^2\setminus \{(0,0)\}$ y $w\in \mathbb{R}^2$ la única geodésica en $C$ que pasa por $G(p)$ con velocidad $(dG)_p(w)\in T_{G(p)}C$:
        \end{enumerate}
    \end{ejercicio}

    \begin{ejercicio}
        Razonar si son verdaderas o falsas las siguientes cuestiones.
        \begin{enumerate}[label=\alph*)]
            \item Si $S_1,S_2$ son superficies regulares orientables isométricas en $\mathbb{R}^3$, entonces $S_1$ es mínima si y sólo si $S_2$ es mínima.
            \item El hiperboloide $\{(x,y,z)\in \mathbb{R}^3 : x^2+y^2 = z^2+1\}$ es isométrico al cilindro $\bb{S}^1\times \mathbb{R}$.
        \end{enumerate}
    \end{ejercicio}

\end{document}
